% Section 16: The six-port whole-cut kernel
\label{sec:six-port-body}

Part~II has so far supplied a semantics, a calculus, and a metatheory.  This section asks
how few ports a carrier needs in order to express that calculus with three readings and
binary cuts, and answers with a counting theorem.  The answer is six.  The section is
careful about what the number means: it is a conditional minimum under a stated contract,
not an absolute truth, and the universality claim that usually accompanies such a number is
replaced by relative initiality.

\subsection{The counting question}

\begin{definition}[Whole-cut carrier]\label{def:whole-cut-carrier}
A \emph{whole-cut carrier} with $m$ domains and $n$ ports consists of
\begin{enumerate}
\item a finite set $D$ of domains;
\item for each domain, a binary cut $\langle L\mid R\rangle$ that is an atomic generator
  rather than a derived construct;
\item a set $P$ of ports, each attached to one domain and one side of its cut;
\item an occurrence-and-multiplicity ledger $M$ recording, for every term, the
  production and consumption of each occurrence;
\item a routing datum $\gamma$, a history ledger $H$, and a residual record $R$.
\end{enumerate}
\end{definition}

\begin{definition}[The contract]\label{def:six-port-contract}
The counting theorem is stated under six clauses.
\begin{enumerate}
\item \textbf{Atomic cuts.}  Each cut is a generator; it is not decomposed into smaller
  structure.
\item \textbf{Binary sides.}  Each cut has exactly two complementary sides.
\item \textbf{Two-sided access.}  Every domain must be readable on both sides of its own
  cut, since a view that could read one side only could not express a crossing.
\item \textbf{Occurrence fidelity.}  No view, braid, collision, projection, or quotient may
  create, drop, rename, or merge an occurrence.
\item \textbf{Production witnesses.}  Copies require production witnesses; collisions
  conserve multiplicity.
\item \textbf{View discipline.}  Every view extends the ledger and commutes with the
  boundary operator: $V_\nu(\partial c)=\partial V_\nu(c)$.
\end{enumerate}
\end{definition}

\subsection{The minimality theorem}

\begin{theorem}[Counting bound]\label{thm:counting-bound}
Under the contract of Definition~\ref{def:six-port-contract}, a triadic whole-cut carrier
with binary cuts satisfies
\[
  n\ \ge\ 2m ,
\]
and hence, for $m=3$, $n_{\min}=6$.
\end{theorem}

\begin{proof}
By clause (3) every domain must be accessible on both sides of its own cut, so each domain
contributes at least two distinct ports; by clause (1) the cut is atomic, so those two ports
cannot be served by sub-structure; by clause (2) the two sides are complementary, so a
single port cannot serve both.  Hence $n\ge 2m$, with equality attained by the carrier
$\bW_6$ of \S\ref{sec:six-port-carrier}.
\end{proof}

\begin{remark}[What the proof does and does not use]
The bound uses clauses (1)--(3) only; clauses (4)--(6) are needed for the carrier to be
usable, not for the count.  This split matters, because it isolates exactly which part of
the contract is under pressure if the number six is ever challenged.
\end{remark}

\begin{warning}[Six is conditional]\label{warn:six-conditional}
Theorem~\ref{thm:counting-bound} says that six is minimal \emph{for this contract}.  Drop
clause (3) and a carrier with three ports becomes available; strengthen clause (1) to
require decomposed cuts and the minimum changes again.  The number six is therefore a
statement about the contract, and reporting it as an absolute minimality of any
three-domain syntax would be an error.
\end{warning}

\subsection{The carrier}\label{sec:six-port-carrier}

\begin{definition}[The six-port carrier]\label{def:w6}
The carrier is
\[
  \bW_6=(\mathsf D,\ C,\ P,\ O,\ M,\ \gamma,\ H,\ R)
\]
consisting of domains, cuts, ports, occurrences, the multiplicity ledger, routing, history,
and the residual record.  Its distinguishing structural features are:
\begin{enumerate}
\item a \emph{split pair}, giving the two sides of each atomic cut;
\item a \emph{projector}, identifying which ports a given view reads;
\item a \emph{factored traversal}, so that a traversal decomposes along the cut rather than
  cutting across it;
\item \emph{collision strata}, recording where two occurrences coincide.
\end{enumerate}
\end{definition}

\begin{proposal}[Line and circle are views]\label{prop:line-circle-views}
In $\bW_6$ the \emph{line} and the \emph{circle} are not neighbouring value constructors but
checked projections of one ledger, each with a declared kernel.  One atomic cut
co-generates a continuation germ and a boundary germ; the line reads the first and the
circle the second.
\end{proposal}

\begin{proof}[Status]
A definitional restructuring, justified by clause (6): both readings are views
$V_\nu$ of the same boundary operator, so neither is primitive.  The claim asserted is the
commutation $V_\nu(\partial c)=\partial V_\nu(c)$, which holds by construction of the two
views and is checked on the finite fixture.
\end{proof}

\subsection{The grammar, the ledgers, and the degree}

\begin{definition}[Whole-cut record]\label{def:whole-cut-record}
A \emph{whole-cut record} consists of port, cut, and cell records together with six ledgers:
\emph{production}, \emph{through}, \emph{cycle}, \emph{hole}, \emph{alternative}, and
\emph{residual}.  The record carries a \emph{degree} $\deg(W)$, and for the minimal carrier
$\deg(\bW_6)=6$.
\end{definition}

\begin{proposition}[Collision conservation]\label{prop:collision-conservation}
In any well-typed whole-cut record the multiplicity ledger is conserved under every rule of
the grammar: no production is created or destroyed, and every collision accounts for
exactly one coincidence of two occurrences.
\end{proposition}

\begin{proof}
Each grammar rule consumes and produces ledger entries in matching multiplicity; the
production ledger records copies explicitly, so a rule that duplicates a term must record a
production, and clause (5) forbids doing so silently.  Induction over derivations in the
grammar.
\end{proof}

\begin{remark}[Distributivity as a recorded witness]
The grammar's distributivity step is accompanied by an explicit witness record with the
ordered hole bindings.  The content is not the equation --- which holds in the intended
model --- but the record of which occurrences of the duplicated term were produced, in
which order.  Without that record the equation would silently license a contraction, which
No-go~\ref{nogo:raw-ledger} and Proposition~\ref{prop:collision-conservation} forbid.
\end{remark}

\subsection{The cell--carrier--view spine}

\begin{proposal}[Spine]\label{prop:cell-carrier-view}
There is a typed spine
\[
  \mathsf{TypedCellSignature}
  \xrightarrow{\ \mathsf{assemble}_\rho\ }
  \mathsf{StratifiedCarrier}
  \xrightarrow{\ V_\nu\ }
  \mathsf{PresentedView}_\nu ,
\]
in which the first arrow assembles cells into a stratified carrier along a declared
stratification $\rho$, and the second is the view of clause (6).  The boundary operators of
the two basic cells are
\[
  \partial\mathsf{InterchangeCell}=Q_4,\qquad
  \partial\mathsf{BraidCell}=M_6 ,
\]
with Coxeter shadows $C_2\times C_2$ and $S_3$ respectively.
\end{proposal}

\begin{proof}[Status]
The boundary identifications are computations in the cell complex: an interchange cell is a
square whose boundary reads as a product of two transpositions, a braid cell a hexagon whose
boundary reads as the three-strand braid relation.  The Coxeter shadows are the standard
reflections of the square and hexagon.
\end{proof}

\begin{proposition}[Envelope count]\label{prop:envelope}
The truncated-octahedral envelope of the stratum satisfies
\[
  TO_{24}=6\,Q_4+8\,M_6 ,
\]
that is, six squares and eight hexagons on twenty-four vertices, and the relation boundary
count is $2m$ in the sense of Theorem~\ref{thm:counting-bound}.
\end{proposition}

\begin{proof}
$6\cdot4+8\cdot6=72$ face--vertex incidences on $24$ vertices with degree $3$:
$24\cdot3=72$.  The consistency of the two counts witnesses that the envelope closes.
\end{proof}

\subsection{The $M_6$ boundary bridge}

\begin{theorem}[Alternating-pairing bridge]\label{thm:m6-bridge}
Let $\kappa$ and $\tau$ be two fixed-point-free involutions on the six-port ledger whose
union is connected.  Then the union of their graphs is a two-regular graph on six vertices,
hence a six-cycle $C_6$, whose edges alternate between the two kinds.  Consequently the
checked incidence data present the open boundary
\[
  \mathsf{bridge}(\bW_6)=M_6^{\partial}\bigl[\kappa\tau\kappa
  \Rightarrow\tau\kappa\tau\bigr],
\]
with no filler supplied.
\end{theorem}

\begin{proof}
A fixed-point-free involution on six elements is a perfect matching, so it contributes three
edges.  Two matchings whose union is connected give a two-regular graph, which is a disjoint
union of cycles; connectivity forces a single six-cycle.  Alternation of edge kinds follows
from the two matchings being distinct, and the displayed presentation of $M_6^{\partial}$
records the two alternating readings of the cycle.  The absence of a filler is
Calibration~\ref{cal:m6-open}.
\end{proof}

\begin{calibration}[The filler is open]\label{cal:m6-open}
The bridge of Theorem~\ref{thm:m6-bridge} presents the boundary $M_6^{\partial}$ as an open
object: the two alternating readings are exhibited, and no 2-cell filling them is
constructed.  Every construction in this material that touches $M_6$ lists the filler as an
open residual.
\end{calibration}

\subsection{Relative initiality instead of universality}

\begin{proposal}[Relative initiality]\label{prop:relative-initiality}
$\bW_6$ is initial in the category of presented models satisfying the contract of
Definition~\ref{def:six-port-contract}: for every such model $N$ there is a morphism
$\bW_6\to N$, unique up to the declared decoration.
\end{proposal}

\begin{proof}[Status]
The statement is a \textsf{Proposal}.  What is available is the counting argument of
Theorem~\ref{thm:counting-bound} together with the view discipline of clause (6), which
supply the candidate morphisms; no uniqueness proof is given here.
\end{proof}

\begin{warning}[Not a universality theorem]\label{warn:not-universal}
Proposition~\ref{prop:relative-initiality} replaces, and does not establish, a claim of
universality.  In particular it says nothing about computation-universality, and it does not
assert that every three-domain syntax factors through $\bW_6$.  The distinction is the
departure point of \S\ref{sec:boundaries}.
\end{warning}

\subsection{Summary}

\begin{enumerate}
\item Six is the minimum under a stated six-clause contract
  (Definitions~\ref{def:whole-cut-carrier}, \ref{def:six-port-contract},
  Theorem~\ref{thm:counting-bound}, Warning~\ref{warn:six-conditional}).
\item In the minimal carrier, line and circle are views of one ledger
  (Definition~\ref{def:w6}, Proposal~\ref{prop:line-circle-views}).
\item The ledgers conserve multiplicity, so distributivity must record its productions
  (Proposition~\ref{prop:collision-conservation}).
\item The spine and envelope counts are explicit
  (Proposal~\ref{prop:cell-carrier-view}, Proposition~\ref{prop:envelope}).
\item The $M_6$ boundary is presented with an open filler, and the kernel's status is
  relative initiality, not universality
  (Theorem~\ref{thm:m6-bridge}, Calibration~\ref{cal:m6-open},
  Proposal~\ref{prop:relative-initiality}, Warning~\ref{warn:not-universal}).
\end{enumerate}
